\chapter{Complexity — Path Hard Regs}
%%% generated by the blueprint pipeline from TCSlib.Complexity.SpaceComplexity.Logspace.PathHardRegs
\section{Overview}

Declarations of the Lean module \texttt{TCSlib.Complexity.SpaceComplexity.Logspace.PathHardRegs}.

\section{Declarations}

\begin{definition}[Input-length register]
\lean{Complexity.PathHard.rY}
\label{Complexity.PathHard.rY}
\leanok
Among fifteen registers, designate register $0$ as the input-length register.
\end{definition}

\begin{definition}[Pointer register]
\lean{Complexity.PathHard.rP}
\label{Complexity.PathHard.rP}
\leanok
Among fifteen registers, designate register $1$ as the pointer register.
\end{definition}

\begin{definition}[Input-value register]
\lean{Complexity.PathHard.rX}
\label{Complexity.PathHard.rX}
\leanok
Among fifteen registers, designate register $2$ as the input-value register.
\end{definition}

\begin{definition}[Temporary register]
\lean{Complexity.PathHard.rT}
\label{Complexity.PathHard.rT}
\leanok
Among fifteen registers, designate register $3$ as the temporary register.
\end{definition}

\begin{definition}[Index c]
\lean{Complexity.PathHard.rC}
\label{Complexity.PathHard.rC}
\leanok
Let $r_C$ denote the element $4$ of the set $\mathbin{15}=\\{0,1,\\ldots,14\\}$, the set of
natural numbers less than $15$.
\end{definition}

\begin{definition}[Index s]
\lean{Complexity.PathHard.rS}
\label{Complexity.PathHard.rS}
\leanok
Let $r_S$ denote the element $5$ of the set $\mathbin{15}=\\{0,1,\\ldots,14\\}$, the set of
natural numbers less than $15$.
\end{definition}

\begin{definition}[Index e]
\lean{Complexity.PathHard.rE}
\label{Complexity.PathHard.rE}
\leanok
Let $r_E$ denote the element $6$ of the set $\mathbin{15}=\\{0,1,\\ldots,14\\}$, the set of
natural numbers less than $15$.
\end{definition}

\begin{definition}[Index n]
\lean{Complexity.PathHard.rN}
\label{Complexity.PathHard.rN}
\leanok
Let $r_N$ denote the element $7$ of the set $\mathbin{15}=\\{0,1,\\ldots,14\\}$, the set of
natural numbers less than $15$.
\end{definition}

\begin{definition}[Index m]
\lean{Complexity.PathHard.rM}
\label{Complexity.PathHard.rM}
\leanok
Let $r_M$ denote the element $8$ of the set $\mathbin{15}=\\{0,1,\\ldots,14\\}$, the set of
natural numbers less than $15$.
\end{definition}

\begin{definition}[Index k]
\lean{Complexity.PathHard.rK}
\label{Complexity.PathHard.rK}
\leanok
Let $r_K$ denote the element $9$ of the set $\mathbin{15}=\\{0,1,\\ldots,14\\}$, the set of
natural numbers less than $15$.
\end{definition}

\begin{definition}[Index q]
\lean{Complexity.PathHard.rQ}
\label{Complexity.PathHard.rQ}
\leanok
Let $r_Q$ denote the element $10$ of the set $\mathbin{15}=\\{0,1,\\ldots,14\\}$, the set of
natural numbers less than $15$.
\end{definition}

\begin{definition}[Index i]
\lean{Complexity.PathHard.rI}
\label{Complexity.PathHard.rI}
\leanok
Let $r_I$ denote the element $11$ of the set $\mathbin{15}=\\{0,1,\\ldots,14\\}$, the set of
natural numbers less than $15$.
\end{definition}

\begin{definition}[Destination register index]
\lean{Complexity.PathHard.rD}
\label{Complexity.PathHard.rD}
\leanok
The destination register is the register with index $12$ among the $15$ registers.
\end{definition}

\begin{definition}[Uppercase register index]
\lean{Complexity.PathHard.rU}
\label{Complexity.PathHard.rU}
\leanok
The uppercase register is the register with index $13$ among the $15$ registers.
\end{definition}

\begin{definition}[Z register index]
\lean{Complexity.PathHard.rZ}
\label{Complexity.PathHard.rZ}
\leanok
The $Z$ register is the register with index $14$ among the $15$ registers.
\end{definition}

\begin{definition}[Machine Configuration Registers]
\lean{Complexity.PathHard.St}
\label{Complexity.PathHard.St}
\leanok
A machine configuration is a tuple of fifteen natural-number registers: the query length
and decoding pointer, the pair-scan pointer, the input length, a scratch register, a
loop counter, a radius, a vertex exponent, registers for $2n$ and $2n^2$, a width cap, a
width-check pointer, an index and remaining offset, decoding and comparison scratch, a
source vertex, and shift scratch.
\end{definition}

\begin{definition}[Register vector representation]
\lean{Complexity.PathHard.St.toV}
\label{Complexity.PathHard.St.toV}
\leanok
\uses{Complexity.PathHard.St}
For a machine configuration, its register vector is the function from the $15$ register
indices to the natural numbers whose entries, in order, are the query-length/decode
pointer, pair-scan pointer, input length, scratch register, loop counter, radius, vertex
exponent, $2n$, $2n^2$, width cap, width-check pointer, index/remaining offset,
decoding/comparison scratch, source vertex, and shift scratch.
\end{definition}

\begin{definition}[Zero configuration]
\lean{Complexity.PathHard.St.zero}
\label{Complexity.PathHard.St.zero}
\leanok
\uses{Complexity.PathHard.St}
The zero configuration is the machine configuration in which all fifteen registers are
zero.
\end{definition}

\begin{lemma}[Zero State Register Vector]
\lean{Complexity.PathHard.St.toV_zero}
\label{Complexity.PathHard.St.toV_zero}
\leanok
\uses{Complexity.PathHard.St.toV, Complexity.PathHard.St.zero}
\difficulty{3}
The register vector of the state whose fifteen registers are all zero is the function
assigning zero to every position.
\end{lemma}

\begin{lemma}[Query Length Coordinate]
\lean{Complexity.PathHard.St.toV_y}
\label{Complexity.PathHard.St.toV_y}
\leanok
\uses{Complexity.PathHard.St, Complexity.PathHard.St.toV, Complexity.PathHard.rY}
\difficulty{1}
For a state $st$ with query-length register $y$, the coordinate of its register vector
at the query-length position is $y$.
\end{lemma}

\begin{lemma}[State vector pair scan pointer]
\lean{Complexity.PathHard.St.toV_p}
\label{Complexity.PathHard.St.toV_p}
\leanok
\uses{Complexity.PathHard.St, Complexity.PathHard.St.toV, Complexity.PathHard.rP}
\difficulty{1}
For every state $st$, the register vector’s coordinate with index $1$ is the state’s
pair-scan pointer $p$.
\end{lemma}

\begin{lemma}[Input Length Coordinate]
\lean{Complexity.PathHard.St.toV_x}
\label{Complexity.PathHard.St.toV_x}
\leanok
\uses{Complexity.PathHard.St, Complexity.PathHard.St.toV, Complexity.PathHard.rX}
\difficulty{1}
For a state $st$ with input-length register $x$, the coordinate of its register vector
at the input-length position is $x$.
\end{lemma}

\begin{lemma}[State vector scratch register]
\lean{Complexity.PathHard.St.toV_t}
\label{Complexity.PathHard.St.toV_t}
\leanok
\uses{Complexity.PathHard.St, Complexity.PathHard.St.toV, Complexity.PathHard.rT}
\difficulty{1}
For every state $st$, the register vector’s coordinate with index $3$ is the state’s
scratch register $t$.
\end{lemma}

\begin{lemma}[Loop-counter coordinate]
\lean{Complexity.PathHard.St.toV_c}
\label{Complexity.PathHard.St.toV_c}
\leanok
\uses{Complexity.PathHard.St, Complexity.PathHard.St.toV, Complexity.PathHard.rC}
\difficulty{1}
For every state record $\mathrm{st}$, the loop-counter coordinate of its fifteen-entry
coordinate vector is the loop-counter value $c$ of $\mathrm{st}$.
\end{lemma}

\begin{lemma}[State vector radius]
\lean{Complexity.PathHard.St.toV_s}
\label{Complexity.PathHard.St.toV_s}
\leanok
\uses{Complexity.PathHard.St, Complexity.PathHard.St.toV, Complexity.PathHard.rS}
\difficulty{1}
For every state $st$, the register vector’s coordinate with index $5$ is the state’s
radius register $s$.
\end{lemma}

\begin{lemma}[Vertex-exponent coordinate]
\lean{Complexity.PathHard.St.toV_e}
\label{Complexity.PathHard.St.toV_e}
\leanok
\uses{Complexity.PathHard.St, Complexity.PathHard.St.toV, Complexity.PathHard.rE}
\difficulty{1}
For every state record $\mathrm{st}$, the vertex-exponent coordinate of its
fifteen-entry coordinate vector is the vertex exponent $e$ of $\mathrm{st}$.
\end{lemma}

\begin{lemma}[State vector vertex count]
\lean{Complexity.PathHard.St.toV_n}
\label{Complexity.PathHard.St.toV_n}
\leanok
\uses{Complexity.PathHard.St, Complexity.PathHard.St.toV, Complexity.PathHard.rN}
\difficulty{1}
For every state $st$, the register vector’s coordinate with index $7$ is the state’s
vertex count register $n$.
\end{lemma}

\begin{lemma}[State vector size register]
\lean{Complexity.PathHard.St.toV_m}
\label{Complexity.PathHard.St.toV_m}
\leanok
\uses{Complexity.PathHard.St, Complexity.PathHard.St.toV, Complexity.PathHard.rM}
\difficulty{1}
For every state $st$, the register vector’s coordinate with index $8$ is the state’s
size register $m$.
\end{lemma}

\begin{lemma}[State vector width cap]
\lean{Complexity.PathHard.St.toV_k}
\label{Complexity.PathHard.St.toV_k}
\leanok
\uses{Complexity.PathHard.St, Complexity.PathHard.St.toV, Complexity.PathHard.rK}
\difficulty{1}
For every state $st$, the register vector’s coordinate with index $9$ is the state’s
width cap $k$.
\end{lemma}

\begin{lemma}[State vector width-check pointer]
\lean{Complexity.PathHard.St.toV_q}
\label{Complexity.PathHard.St.toV_q}
\leanok
\uses{Complexity.PathHard.St, Complexity.PathHard.St.toV, Complexity.PathHard.rQ}
\difficulty{1}
For every state $st$, the register vector’s coordinate with index $10$ is the state’s
width-check pointer $q$.
\end{lemma}

\begin{lemma}[State vector index register]
\lean{Complexity.PathHard.St.toV_i}
\label{Complexity.PathHard.St.toV_i}
\leanok
\uses{Complexity.PathHard.St, Complexity.PathHard.St.toV, Complexity.PathHard.rI}
\difficulty{1}
For every state $st$, the register vector’s coordinate with index $11$ is the state’s
index register $i$.
\end{lemma}

\begin{lemma}[Scratch-coordinate value]
\lean{Complexity.PathHard.St.toV_d}
\label{Complexity.PathHard.St.toV_d}
\leanok
\uses{Complexity.PathHard.St, Complexity.PathHard.St.toV, Complexity.PathHard.rD}
\difficulty{1}
For every state record $\mathrm{st}$, the decoding-and-comparison scratch coordinate of
its fifteen-entry coordinate vector is the scratch value $d$ of $\mathrm{st}$.
\end{lemma}

\begin{lemma}[Source Register Coordinate]
\lean{Complexity.PathHard.St.toV_u}
\label{Complexity.PathHard.St.toV_u}
\leanok
\uses{Complexity.PathHard.St, Complexity.PathHard.St.toV, Complexity.PathHard.rU}
\difficulty{1}
For a state $st$ with source-vertex register $u$, the coordinate of its register vector
at the source-vertex position is $u$.
\end{lemma}

\begin{lemma}[Shift Scratch Coordinate]
\lean{Complexity.PathHard.St.toV_z}
\label{Complexity.PathHard.St.toV_z}
\leanok
\uses{Complexity.PathHard.St, Complexity.PathHard.St.toV, Complexity.PathHard.rZ}
\difficulty{1}
For a state $st$ with shift-scratch register $z$, the coordinate of its register vector
at the shift-scratch position is $z$.
\end{lemma}

\begin{lemma}[Updating the query-length register]
\lean{Complexity.PathHard.St.upd_y}
\label{Complexity.PathHard.St.upd_y}
\leanok
\uses{Complexity.PathHard.St, Complexity.PathHard.St.toV, Complexity.PathHard.rY}
\difficulty{4}
For a state $s$ with fifteen natural-number registers, let its register vector list the
values of $y,p,x,t,c,s,e,n,m,k,q,i,d,u,z$ in that order. Replacing coordinate $0$ of
this vector by $v$ gives the register vector of the state obtained by replacing its
query-length register $y$ by $v$.
\end{lemma}

\begin{lemma}[Updating the Pair-Scan Pointer]
\lean{Complexity.PathHard.St.upd_p}
\label{Complexity.PathHard.St.upd_p}
\leanok
\uses{Complexity.PathHard.St, Complexity.PathHard.St.toV, Complexity.PathHard.rP}
\difficulty{4}
Let a state consist of fifteen natural-number registers, ordered as the query length,
pair-scan pointer, input length, scratch register, loop counter, radius, vertex
exponent, twice the vertex count, twice its square, width cap, width-check pointer,
index, decoding scratch, source vertex, and shift scratch. Replacing the second
coordinate of its register vector by a natural number $v$ gives the register vector of
the state obtained by setting its pair-scan pointer to $v$.
\end{lemma}

\begin{lemma}[Updating the input-length register]
\lean{Complexity.PathHard.St.upd_x}
\label{Complexity.PathHard.St.upd_x}
\leanok
\uses{Complexity.PathHard.St, Complexity.PathHard.St.toV, Complexity.PathHard.rX}
\difficulty{4}
For a state $s$ with fifteen natural-number registers, let its register vector list the
values of $y,p,x,t,c,s,e,n,m,k,q,i,d,u,z$ in that order. Replacing coordinate $2$ of
this vector by $v$ gives the register vector of the state obtained by replacing its
input-length register $x$ by $v$.
\end{lemma}

\begin{lemma}[Updating the Scratch Register]
\lean{Complexity.PathHard.St.upd_t}
\label{Complexity.PathHard.St.upd_t}
\leanok
\uses{Complexity.PathHard.St, Complexity.PathHard.St.toV, Complexity.PathHard.rT}
\difficulty{4}
Let a state consist of fifteen natural-number registers, ordered as the query length,
pair-scan pointer, input length, scratch register, loop counter, radius, vertex
exponent, twice the vertex count, twice its square, width cap, width-check pointer,
index, decoding scratch, source vertex, and shift scratch. Replacing the fourth
coordinate of its register vector by a natural number $v$ gives the register vector of
the state obtained by setting its scratch register to $v$.
\end{lemma}

\begin{lemma}[Loop Counter Update]
\lean{Complexity.PathHard.St.upd_c}
\label{Complexity.PathHard.St.upd_c}
\leanok
\uses{Complexity.PathHard.St, Complexity.PathHard.St.toV, Complexity.PathHard.rC}
\difficulty{4}
For any state $st$ and natural number $v$, replacing the loop-counter coordinate of its
register vector by $v$ gives the register vector of the state obtained from $st$ by
setting its loop counter to $v$.
\end{lemma}

\begin{lemma}[Updating the Radius Coordinate]
\lean{Complexity.PathHard.St.upd_s}
\label{Complexity.PathHard.St.upd_s}
\leanok
\uses{Complexity.PathHard.St, Complexity.PathHard.St.toV, Complexity.PathHard.rS}
\difficulty{4}
Let a state consist of fifteen natural-number registers, ordered as the query length,
pair-scan pointer, input length, scratch register, loop counter, radius, vertex
exponent, twice the vertex count, twice its square, width cap, width-check pointer,
index, decoding scratch, source vertex, and shift scratch. Replacing the sixth
coordinate of its register vector by a natural number $v$ gives the register vector of
the state obtained by setting its radius register to $v$.
\end{lemma}

\begin{lemma}[Vertex Exponent Update]
\lean{Complexity.PathHard.St.upd_e}
\label{Complexity.PathHard.St.upd_e}
\leanok
\uses{Complexity.PathHard.St, Complexity.PathHard.St.toV, Complexity.PathHard.rE}
\difficulty{4}
For any state $st$ and natural number $v$, replacing the vertex-exponent coordinate of
its register vector by $v$ gives the register vector of the state obtained from $st$ by
setting its vertex exponent to $v$.
\end{lemma}

\begin{lemma}[Updating the Vertex-Count Coordinate]
\lean{Complexity.PathHard.St.upd_n}
\label{Complexity.PathHard.St.upd_n}
\leanok
\uses{Complexity.PathHard.St, Complexity.PathHard.St.toV, Complexity.PathHard.rN}
\difficulty{4}
Let a state consist of fifteen natural-number registers, ordered as the query length,
pair-scan pointer, input length, scratch register, loop counter, radius, vertex
exponent, twice the vertex count, twice its square, width cap, width-check pointer,
index, decoding scratch, source vertex, and shift scratch. Replacing the eighth
coordinate of its register vector by a natural number $v$ gives the register vector of
the state obtained by setting its register for twice the vertex count to $v$.
\end{lemma}

\begin{lemma}[Updating the Square-Count Coordinate]
\lean{Complexity.PathHard.St.upd_m}
\label{Complexity.PathHard.St.upd_m}
\leanok
\uses{Complexity.PathHard.St, Complexity.PathHard.St.toV, Complexity.PathHard.rM}
\difficulty{3}
Let a state consist of fifteen natural-number registers, ordered as the query length,
pair-scan pointer, input length, scratch register, loop counter, radius, vertex
exponent, twice the vertex count, twice its square, width cap, width-check pointer,
index, decoding scratch, source vertex, and shift scratch. Replacing the ninth
coordinate of its register vector by a natural number $v$ gives the register vector of
the state obtained by setting its register for twice the square of the vertex count to
$v$.
\end{lemma}

\begin{lemma}[Updating the Width-Cap Coordinate]
\lean{Complexity.PathHard.St.upd_k}
\label{Complexity.PathHard.St.upd_k}
\leanok
\uses{Complexity.PathHard.St, Complexity.PathHard.St.toV, Complexity.PathHard.rK}
\difficulty{3}
Let a state consist of fifteen natural-number registers, ordered as the query length,
pair-scan pointer, input length, scratch register, loop counter, radius, vertex
exponent, twice the vertex count, twice its square, width cap, width-check pointer,
index, decoding scratch, source vertex, and shift scratch. Replacing the tenth
coordinate of its register vector by a natural number $v$ gives the register vector of
the state obtained by setting its width-cap register to $v$.
\end{lemma}

\begin{lemma}[Updating the Width-Check Pointer]
\lean{Complexity.PathHard.St.upd_q}
\label{Complexity.PathHard.St.upd_q}
\leanok
\uses{Complexity.PathHard.St, Complexity.PathHard.St.toV, Complexity.PathHard.rQ}
\difficulty{3}
Let a state consist of fifteen natural-number registers, ordered as the query length,
pair-scan pointer, input length, scratch register, loop counter, radius, vertex
exponent, twice the vertex count, twice its square, width cap, width-check pointer,
index, decoding scratch, source vertex, and shift scratch. Replacing the eleventh
coordinate of its register vector by a natural number $v$ gives the register vector of
the state obtained by setting its width-check pointer to $v$.
\end{lemma}

\begin{lemma}[Updating the Index Coordinate]
\lean{Complexity.PathHard.St.upd_i}
\label{Complexity.PathHard.St.upd_i}
\leanok
\uses{Complexity.PathHard.St, Complexity.PathHard.St.toV, Complexity.PathHard.rI}
\difficulty{3}
Let a state consist of fifteen natural-number registers, ordered as the query length,
pair-scan pointer, input length, scratch register, loop counter, radius, vertex
exponent, twice the vertex count, twice its square, width cap, width-check pointer,
index, decoding scratch, source vertex, and shift scratch. Replacing the twelfth
coordinate of its register vector by a natural number $v$ gives the register vector of
the state obtained by setting its index register to $v$.
\end{lemma}

\begin{lemma}[Comparison Scratch Update]
\lean{Complexity.PathHard.St.upd_d}
\label{Complexity.PathHard.St.upd_d}
\leanok
\uses{Complexity.PathHard.St, Complexity.PathHard.St.toV, Complexity.PathHard.rD}
\difficulty{3}
For any state $st$ and natural number $v$, replacing the decoding-and-comparison scratch
coordinate of its register vector by $v$ gives the register vector of the state obtained
from $st$ by setting that scratch register to $v$.
\end{lemma}

\begin{lemma}[Updating the source register]
\lean{Complexity.PathHard.St.upd_u}
\label{Complexity.PathHard.St.upd_u}
\leanok
\uses{Complexity.PathHard.St, Complexity.PathHard.St.toV, Complexity.PathHard.rU}
\difficulty{3}
For a state $s$ with fifteen natural-number registers, let its register vector list the
values of $y,p,x,t,c,s,e,n,m,k,q,i,d,u,z$ in that order. Replacing coordinate $13$ of
this vector by $v$ gives the register vector of the state obtained by replacing its
source register $u$ by $v$.
\end{lemma}

\begin{lemma}[Updating the shift register]
\lean{Complexity.PathHard.St.upd_z}
\label{Complexity.PathHard.St.upd_z}
\leanok
\uses{Complexity.PathHard.St, Complexity.PathHard.St.toV, Complexity.PathHard.rZ}
\difficulty{4}
For a state $s$ with fifteen natural-number registers, let its register vector list the
values of $y,p,x,t,c,s,e,n,m,k,q,i,d,u,z$ in that order. Replacing coordinate $14$ of
this vector by $v$ gives the register vector of the state obtained by replacing its
shift register $z$ by $v$.
\end{lemma}

\begin{definition}[Register bound]
\lean{Complexity.PathHard.St.Bd}
\label{Complexity.PathHard.St.Bd}
\leanok
\uses{Complexity.PathHard.St}
A machine configuration is bounded by a natural number $B$ if each of its fifteen
registers is at most $B$.
\end{definition}

\begin{lemma}[State bounds imply coordinate bounds]
\lean{Complexity.PathHard.St.Bd.toV}
\label{Complexity.PathHard.St.Bd.toV}
\leanok
\uses{Complexity.PathHard.St, Complexity.PathHard.St.Bd, Complexity.PathHard.St.toV}
\difficulty{4}
Let $B$ be a natural number and let $\mathrm{st}$ be a state record with fifteen
natural-number coordinates, whose coordinate vector lists those coordinates in order. If
each coordinate of $\mathrm{st}$ is at most $B$, then every entry of its coordinate
vector is at most $B$.
\end{lemma}

\begin{lemma}[Coordinate bounds imply state bounds]
\lean{Complexity.PathHard.St.Bd.of_toV}
\label{Complexity.PathHard.St.Bd.of_toV}
\leanok
\uses{Complexity.PathHard.St, Complexity.PathHard.St.Bd, Complexity.PathHard.St.toV, Complexity.PathHard.rC, Complexity.PathHard.rD, Complexity.PathHard.rE, Complexity.PathHard.rI, Complexity.PathHard.rK, Complexity.PathHard.rM, Complexity.PathHard.rN, Complexity.PathHard.rP, Complexity.PathHard.rQ, Complexity.PathHard.rS, Complexity.PathHard.rT, Complexity.PathHard.rU, Complexity.PathHard.rX, Complexity.PathHard.rY, Complexity.PathHard.rZ}
\difficulty{3}
Let $B$ be a natural number and let $\mathrm{st}$ be a state record with fifteen
natural-number coordinates, whose coordinate vector lists those coordinates in order. If
every entry of this vector is at most $B$, then each of the fifteen coordinates of
$\mathrm{st}$ is at most $B$.
\end{lemma}

\begin{lemma}[Reconstructing a state from two coordinates]
\lean{Complexity.PathHard.St.eq_qk}
\label{Complexity.PathHard.St.eq_qk}
\leanok
\uses{Complexity.PathHard.St, Complexity.PathHard.St.toV, Complexity.PathHard.St.toV_c, Complexity.PathHard.St.toV_d, Complexity.PathHard.St.toV_e, Complexity.PathHard.St.toV_i, Complexity.PathHard.St.toV_m, Complexity.PathHard.St.toV_n, Complexity.PathHard.St.toV_p, Complexity.PathHard.St.toV_s, Complexity.PathHard.St.toV_t, Complexity.PathHard.St.toV_u, Complexity.PathHard.St.toV_x, Complexity.PathHard.St.toV_y, Complexity.PathHard.St.toV_z, Complexity.PathHard.rC, Complexity.PathHard.rD, Complexity.PathHard.rE, Complexity.PathHard.rI, Complexity.PathHard.rK, Complexity.PathHard.rM, Complexity.PathHard.rN, Complexity.PathHard.rP, Complexity.PathHard.rQ, Complexity.PathHard.rS, Complexity.PathHard.rT, Complexity.PathHard.rU, Complexity.PathHard.rX, Complexity.PathHard.rY, Complexity.PathHard.rZ}
\difficulty{5}
Let $\mathrm{st}$ be a state record with fifteen natural-number coordinates, and let $v$
be another fifteen-entry vector of natural numbers. Suppose that $v$ agrees with the
coordinate vector of $\mathrm{st}$ at every coordinate other than the width-cap
coordinate and the pointer-of-width-check coordinate. Then $v$ is the coordinate vector
of the state obtained from $\mathrm{st}$ by replacing those two coordinates with the
corresponding entries of $v$.
\end{lemma}

\begin{lemma}[Reconstructing a state from three coordinates]
\lean{Complexity.PathHard.St.eq_yid}
\label{Complexity.PathHard.St.eq_yid}
\leanok
\uses{Complexity.PathHard.St, Complexity.PathHard.St.toV, Complexity.PathHard.St.toV_c, Complexity.PathHard.St.toV_e, Complexity.PathHard.St.toV_k, Complexity.PathHard.St.toV_m, Complexity.PathHard.St.toV_n, Complexity.PathHard.St.toV_p, Complexity.PathHard.St.toV_q, Complexity.PathHard.St.toV_s, Complexity.PathHard.St.toV_t, Complexity.PathHard.St.toV_u, Complexity.PathHard.St.toV_x, Complexity.PathHard.St.toV_z, Complexity.PathHard.rC, Complexity.PathHard.rD, Complexity.PathHard.rE, Complexity.PathHard.rI, Complexity.PathHard.rK, Complexity.PathHard.rM, Complexity.PathHard.rN, Complexity.PathHard.rP, Complexity.PathHard.rQ, Complexity.PathHard.rS, Complexity.PathHard.rT, Complexity.PathHard.rU, Complexity.PathHard.rX, Complexity.PathHard.rY, Complexity.PathHard.rZ}
\difficulty{5}
Let $\mathrm{st}$ be a state record with fifteen natural-number coordinates, and let $v$
be another fifteen-entry vector of natural numbers. Suppose that $v$ agrees with the
coordinate vector of $\mathrm{st}$ at every coordinate other than the query-length,
index, and decoding-and-comparison-scratch coordinates. Then $v$ is the coordinate
vector of the state obtained from $\mathrm{st}$ by replacing those three coordinates
with the corresponding entries of $v$.
\end{lemma}

\begin{definition}[Machine-associated languages]
\lean{Complexity.PathHard.As}
\label{Complexity.PathHard.As}
\leanok
\uses{Complexity.Codec.acceptLang, Complexity.Codec.initLang, Complexity.Codec.stepLang, Complexity.dropIdx, Complexity.inputBitLang, Complexity.inputLenLang, Turing.FinNDTM}
Given a finite nondeterministic Turing machine $N$ over the Boolean alphabet, define a
function assigning a language over Boolean strings to each of the five indices. In
order, these languages consist of strings that are input-bit encodings, input-length
encodings, encodings whose decoded string describes an initial configuration of $N$,
encodings whose decoded string describes an accepting configuration of $N$, and
encodings whose decoded string describes a valid step of $N$.
\end{definition}

\begin{lemma}[Logarithmic-space membership of the encoded languages]
\lean{Complexity.PathHard.As_mem}
\label{Complexity.PathHard.As_mem}
\leanok
\uses{Complexity.Codec.acceptLang_mem, Complexity.Codec.initLang_mem, Complexity.Codec.stepLang_mem, Complexity.LOGSPACE, Complexity.PathHard.As, Complexity.implicitlyLogspaceComputable_dropIdx, Complexity.inputBitLang_mem, Complexity.inputLenLang_mem, Complexity.mem_LOGSPACE_preimage, Turing.FinNDTM}
\difficulty{4}
Let $N$ be a finite-state nondeterministic Turing machine with Boolean input alphabet.
Each of the five languages in the family associated with $N$—encoding input bits, input
lengths, initial configurations, accepting configurations, and single computation
steps—is decidable by a deterministic Turing machine using at most a constant multiple
of $loor{\log_2 n}+1$ space on inputs of length $n$.
\end{lemma}

\begin{definition}[Five associated language indicators]
\lean{Complexity.PathHard.orc}
\label{Complexity.PathHard.orc}
\leanok
\uses{Complexity.PathHard.As, Turing.FinNDTM, Turing.MultiTapeTM.indicator}
For a finite nondeterministic Turing machine $N$ over the Boolean alphabet, define a
function on indices $j\in\{0,1,2,3,4\}$ and Boolean words $V$ by returning $1$ exactly
when $V$ belongs to the language associated with $N$ and $j$, and $0$ otherwise.
\end{definition}

\begin{lemma}[First Oracle Component]
\lean{Complexity.PathHard.orc_bit}
\label{Complexity.PathHard.orc_bit}
\leanok
\uses{Complexity.PathHard.As, Complexity.PathHard.orc, Complexity.indicator_indexLang_pairEncode, Complexity.inputBitLang, Turing.FinNDTM, Turing.pairEncode}
\difficulty{4}
For every finite nondeterministic Turing machine $N$, bit string $y$, and natural number
$p$, the first oracle component, evaluated on the pairing of $y$ with the binary
representation of $p$, returns the bit of $y$ at position $p$, with positions counted
from $0$ and out-of-range positions interpreted as false.
\end{lemma}

\begin{lemma}[Input-Length Oracle Component]
\lean{Complexity.PathHard.orc_len}
\label{Complexity.PathHard.orc_len}
\leanok
\uses{Complexity.PathHard.As, Complexity.PathHard.orc, Complexity.indicator_indexLang_pairEncode, Complexity.inputLenLang, Turing.FinNDTM, Turing.pairEncode}
\difficulty{3}
For every finite nondeterministic Turing machine $N$, bit string $y$, and natural number
$p$, the second oracle component on the pairing of $y$ with the binary representation of
$p$ returns true exactly when $p<|y|$.
\end{lemma}

\begin{lemma}[Initial Configuration Oracle Component]
\lean{Complexity.PathHard.orc_init}
\label{Complexity.PathHard.orc_init}
\leanok
\uses{Complexity.Codec.initLang, Complexity.LogProg.indicator_eq_decide, Complexity.PathHard.As, Complexity.PathHard.orc, Complexity.dropIdx_pairEncode, Turing.FinNDTM, Turing.pairEncode}
\difficulty{3}
For every finite nondeterministic Turing machine $N$, bit strings $x$ and $\beta$, and
natural numbers $v$ and $s$, the third oracle component on the encoded input consisting
of $x$, $\beta$, the binary representation of $v$, and the binary representation of $s$
returns true exactly when the pairing of $x$ with the two binary representations of $v$
and $s$ belongs to the initial-configuration language of $N$.
\end{lemma}

\begin{lemma}[Acceptance oracle on encoded queries]
\lean{Complexity.PathHard.orc_acc}
\label{Complexity.PathHard.orc_acc}
\leanok
\uses{Complexity.Codec.acceptLang, Complexity.LogProg.indicator_eq_decide, Complexity.PathHard.As, Complexity.PathHard.orc, Complexity.dropIdx_pairEncode, Turing.FinNDTM, Turing.pairEncode}
\difficulty{3}
Let $N$ be a finite nondeterministic Turing machine with Boolean tape symbols, and let
$x,\beta$ be bit strings and $u,s$ nonnegative integers. The oracle's answer on the
encoding of the pair whose first component encodes $(x,\beta)$ and whose second
component encodes $(u,s)$ is true exactly when the encoding of $(x,u,s)$ belongs to the
acceptance language of $N$.
\end{lemma}

\begin{lemma}[Transition Oracle Component]
\lean{Complexity.PathHard.orc_step}
\label{Complexity.PathHard.orc_step}
\leanok
\uses{Complexity.Codec.StepRel, Complexity.LogProg.indicator_eq_decide, Complexity.PathHard.As, Complexity.PathHard.orc, Complexity.dropIdx_pairEncode, Turing.FinNDTM, Turing.pairEncode}
\difficulty{3}
For every finite nondeterministic Turing machine $N$, bit strings $x$ and $\beta$, and
natural numbers $u,v,s$, the fifth oracle component on the encoded input consisting of
$x$, $\beta$, and the binary representations of $u,v,s$ returns true exactly when there
is a valid one-step transition of $N$ on input $x$ from a configuration of code $u$ to a
configuration of code $v$, with both configurations confined to the radius-$s$ work-tape
window.
\end{lemma}
